\begin{definition}
Let $n\ge2$, $p>1$, $\varepsilon,\gamma_0,\gamma_1\in\R$, let
$\Omega\subseteq\R^n$ be open and let $A\colon\R^n\to\R^n$.  The predicate
$\mathrm{CitedRegularityEpsPos}(n,p,\varepsilon,\gamma_0,\gamma_1,\Omega,A)$ is the
cited result (C1) of the paper's Note for the formalization
([LU, chapter 4], [DB], [GT, chapter 15]):

\emph{if $\varepsilon>0$ and $A$ satisfies (Q1) of the thesis and (Q2) of the thesis with
parameters $p,\varepsilon,\gamma_0,\gamma_1$ on $\Omega$
(\noderef{EllipticStructureQ12}), then every weak solution
$u\in W^{1,p}_{\mathrm{loc}}(\Omega)$ of $\operatorname{div}A(\nabla u)=0$
(\noderef{IsWeakSolutionOn}) with weak gradient $g$ satisfies
$u\in W^{2,2}_{\mathrm{loc}}(\Omega)$ and $\nabla u\in C(\Omega)$.}

Written out, the conclusion is: $u$ lies in $W^{1,2}_{\mathrm{loc}}(\Omega)$
with weak gradient $g$ (\noderef{MemW1pLocOn}); and there is a field $g'$ with
\begin{enumerate}
\item $g'$ continuous on $\Omega$;
\item $g=g'$ almost everywhere on $\Omega$; and
\item each component $g'^i$ in $W^{1,2}_{\mathrm{loc}}(\Omega)$ with some weak
  gradient (\noderef{MemW1pLocOn}).
\end{enumerate}
The conclusion $\nabla u\in C(\Omega)$ gives in
particular $\nabla u\in L^\infty_{\mathrm{loc}}(\Omega)$ in the same sense, that
of the continuous representative.

\medskip
\noindent\textbf{Why ``$\nabla u\in C(\Omega)$'' is a continuous
representative and not continuity of the given $g$.}  A weak gradient is
determined only almost everywhere: every other clause in sight --- the
weak-solution identity of \noderef{IsWeakSolutionOn}, the weak-gradient identity
and the $L^q$ clauses of \noderef{MemW1pLocOn} --- is unchanged when $g$ is
redefined on a set of measure zero.  Continuity of $g$ is not.  So a version of
this cited result that asserted $\mathrm{ContinuousOn}\;g\;\Omega$ for the
\emph{given} $g$ would be a false $\mathrm{Prop}$: from any weak solution one
redefines $g$ at a single point of $\Omega$, every other clause survives, and
that one fails --- so no $A$ admitting a weak solution would satisfy the
predicate, and every statement assuming it universally would be vacuously true.
The Note at l.~258 forbids exactly this: the cited results ``must not be vacuous
(each must be a true statement for the operators to which it is applied)''.
Clauses (1)--(3) above are the faithful reading of the Note's own wording at
l.~260, ``$\nabla u\in C(\Omega)$'': for an almost-everywhere class, lying in
$C(\Omega)$ means having a representative in $C(\Omega)$.  Clause~(3) is stated
for $g'$ rather than for $g$; this is no change of content, since
$\mathrm{MemW1pLocOn}$ is itself invariant under modification on a null set, and
it keeps every clause about $g'$ in one place.  The same interface decision was
taken for boundary values (pm-0004) and is the shape
\noderef{CitedCampanatoEmbedding} already uses, (C5)'s own wording at l.~266
being ``agrees a.e.\ with a function in $C^{0,\alpha}$''.  As the Note records,
the thesis says ``$C^2$'' at this point, which under (Q1) of the thesis--(Q2) of the thesis
alone is stronger than the literature; the statement above is the weaker one
from the literature, so that the arguments of \S3 must be run with the
differentiated equation holding only in the weak sense.  This result is used as
an explicit hypothesis where it is needed and is never introduced as an axiom.

Under the reading of (Q1) of the thesis fixed in \noderef{EllipticStructureQ12} ---
``$A$ is differentiable with respect to $h$'' formalized as ``$A$ is of class
$C^1$ in $h$ on the set where the weight of (Q1) of the thesis is finite'', which at
the $\varepsilon>0$ of the present statement is all of $\R^n$ --- the statement
above is exactly the literature statement:
[LU, chapter 4], [DB] and [GT, chapter 15] prove it for operators with a
continuous $D_hA$.  Under the merely-differentiable reading this predicate would
assert the same conclusion for a strictly wider class of operators than the
literature covers, which the Note forbids for a cited result; the $C^1$ clause
of \noderef{EllipticStructureQ12} is what keeps it non-vacuous and not stronger
than its sources.
\end{definition}
